\documentclass[reprint,amsmath,amssymb,aps,pra,superscriptaddress]{revtex4-2}
\usepackage{graphicx}
\usepackage{dcolumn}
\usepackage{bm}
\usepackage{hyperref}
\usepackage{physics}
\usepackage{listings}
\usepackage{xcolor}
\usepackage{multirow}
\usepackage{booktabs}
\hypersetup{colorlinks=true,linkcolor=blue,citecolor=blue,urlcolor=blue}
\lstset{backgroundcolor=\color{gray!10},basicstyle=\ttfamily\small,
breaklines=true,keywordstyle=\color{blue},
commentstyle=\color{green!50!black},stringstyle=\color{red},
frame=single,language=Python}
\begin{document}
\title{Dual-Threshold Quantum Decoherence Model (DTQEM):\\
First-Principles Derivation of the v18.0-C Crossover Framework\\
and Statistical Detection Limits}
\author{Reddouane Berramdane}
\email{reddoma@gmail.com}
\affiliation{Independent Researcher, Oujda, Morocco}
\date{\today}

\begin{abstract}
Characterizing quantum decoherence in open quantum systems is essential for
developing robust quantum technologies. In this work we present the physical
and statistical foundations of the Dual-Threshold Quantum Decoherence Model
(DTQEM) version 18.0-C. Unlike previous decoupled formulations, the v18.0-C
framework unifies path-measurement-induced decoherence and thermal dephasing
into a single expression:
$C = C_0\exp(-a I_{\text{path}} - b\,\Delta T/T_{\text{ref}}
- c\,I_{\text{path}}\,\Delta T/T_{\text{ref}})$,
accounting for cross-correlated environmental modes ($\rho>0$) through a
crossover coefficient $c$. The historical model v17.0-C is the
independent-bath limit ($c=0$). We provide first-principles derivations using
the spin-boson framework and Fermi's Golden Rule, deriving $a(\omega_c,T)$
from a Lorentzian spectral density and $b(T)$ from an Ohmic bath. An
operational bridge maps $I_{\text{path}}$ to cavity photon occupancy $\bar n$
via the Englert--Greenberger relation. AICc analysis establishes a strict
statistical threshold: $N\ge36$ samples are required to reliably detect
$c>0$. The framework is extended to predict Entanglement Sudden Death (ESD),
yielding a closed-form $t_{\text{ESD}}$. We further present a theoretical
derivation linking the thermal coefficient $b$ to the molecular bond force
constant $k$ via the Morse potential and Arrhenius kinetics, opening an
interdisciplinary bridge to physical chemistry. All results are fully
reproducible through the open-source code archived at
Zenodo~(DOI:\,\href{https://doi.org/10.5281/zenodo.20516742}{10.5281/zenodo.20516742}).
\end{abstract}
\maketitle

%-------------------------------------------------------------------
\section{Introduction}
%-------------------------------------------------------------------
Maintaining quantum coherence in the presence of environmental noise remains
a central challenge of quantum information processing. Decoherence is driven
by two distinct mechanisms: active extraction of path information by a
detector, and passive thermal dephasing~\cite{breuer2002}. These channels
have historically been treated as independent. The complementarity principle,
framed through the Englert--Greenberger--Yasin (EGY)
inequality~\cite{englert1996}, links path information to fringe visibility,
while thermal dephasing is captured by the spin-boson or Caldeira--Leggett
frameworks~\cite{weiss2012}. In realistic solid-state or
interferometric setups, however, the two baths share physical modes---e.g.,
lattice vibrations acting simultaneously as displacement and thermal baths---
making the independent-bath approximation incomplete.

A useful conceptual image is provided by the hydrogen spectrum. In an isolated
quantum system, discrete energy levels produce sharp spectral lines that serve
as the system's quantum fingerprint. Rising temperature broadens these lines
through collisions and environmental coupling, introducing uncontrolled degrees
of freedom. Decoherence can thus be understood as the gradual loss of a shared
quantum ``language'' between subsystems: a pure, ordered state evolves into a
mixed statistical ensemble with reduced coherence and weakened entanglement.
This picture is consistent with the DTQEM description, in which environmental
correlations and temperature jointly govern coherence decay.

To remedy the independent-bath limitation, we present
\textbf{DTQEM v18.0-C}, which introduces a crossover coupling term
$c\cdot I_{\text{path}}\cdot\Delta T/T_{\text{ref}}$ derived from a
second-order joint cumulant expansion. We justify its structure through
first-principles derivations, provide a fully open-source numerical
implementation, determine the AICc-based statistical detection threshold
for $c$, derive analytical ESD predictions, and propose an
interdisciplinary bridge to physical chemistry via the Morse potential.

%-------------------------------------------------------------------
\section{The Generalized DTQEM Model (v18.0-C)}
%-------------------------------------------------------------------
\subsection{Model Formulation}
Under cross-correlated environments ($\rho>0$), the coherence $C$ is
\begin{equation}
C(I_{\text{path}},T)=C_0\exp\!\left(
-a\,I_{\text{path}}-b\,\frac{\Delta T}{T_{\text{ref}}}
-c\,I_{\text{path}}\,\frac{\Delta T}{T_{\text{ref}}}\right),
\label{eq:v18}
\end{equation}
with $I_{\text{path}}\in[0,1]$, $T_{\text{ref}}=300\,\mathrm{K}$,
$\Delta T=T-T_{\text{ref}}$. Setting $c=0$ recovers \textbf{v17.0-C}.

\subsection{Calibration}
The model was calibrated over $I_{\text{path}}\in[0,1]$ and
$T\in[300,550]\,\mathrm{K}$ using numerical data; see Tables~\ref{tab:parameters}
and~\ref{tab:performance}.

\begin{table}[h]
\caption{Calibrated parameters for DTQEM v18.0-C.}
\label{tab:parameters}
\begin{ruledtabular}
\begin{tabular}{lccc}
Parameter & Value & Std.\ error & Meaning\\\hline
$C_0$ & 0.3675 & $\pm0.0052$ & Baseline coherence\\
$a$ & 1.6968 & $\pm0.0245$ & Path decoherence rate\\
$b$ & 0.8055 & $\pm0.0148$ & Thermal dephasing rate\\
$c$ & 0.5000 & $\pm0.0201$ & Crossover coupling\\
$T_{\text{ref}}$ & 300\,K & fixed & Reference temperature\\
\end{tabular}
\end{ruledtabular}
\end{table}

\begin{table}[h]
\caption{Statistical performance: v18.0-C vs.\ v17.0-C.}
\label{tab:performance}
\begin{ruledtabular}
\begin{tabular}{lccc}
Model & $R^2$ & LOOCV $R^2$ & RMSE\\\hline
v17.0-C & 0.9679 & 0.9356 & 0.0202\\
v18.0-C & 0.9982 & 0.9814 & 0.0045\\
\end{tabular}
\end{ruledtabular}
\end{table}

%-------------------------------------------------------------------
\section{First-Principles Derivations}
%-------------------------------------------------------------------
\subsection{Path-Decoherence Coefficient $a(\omega_c,T)$}
The path-detection channel is modelled as a two-level system coupled to a
bosonic bath~\cite{weiss2012}:
\begin{equation}
H=\frac{\hbar\omega_0}{2}\sigma_z
+\sum_k\hbar\omega_k b_k^\dagger b_k
+\sigma_z\sum_k g_k(b_k^\dagger+b_k).
\end{equation}
Fermi's Golden Rule applied to the Lorentzian spectral density
$J_{\text{det}}(\omega)=\frac{\eta}{\pi}
\frac{\gamma^2}{(\omega-\omega_{\text{path}})^2+\gamma^2}$
gives
\begin{equation}
a(\omega_c,T)=\frac{\eta}{\gamma}
\frac{\coth\!\left(\frac{\hbar\omega_c}{2k_BT}\right)}
{1+\left(\frac{\omega_c-\omega_{\text{path}}}{\gamma}\right)^2}.
\label{eq:a_micro}
\end{equation}

\begin{figure}[t]
\centering
\includegraphics[width=0.85\linewidth]{figures/figure1.png}
\caption{Spectral line shape $a(\omega_c)$; peak at
$\omega_{\text{path}}=5.0\,\mathrm{GHz}$, $\gamma=1.0\,\mathrm{GHz}$.}
\label{fig:lorentzian}
\end{figure}

\subsection{Thermal Dephasing Coefficient $b(T)$}
For an Ohmic bath~\cite{breuer2002} in the high-$T$ limit:
\begin{equation}
b(T)=\left(\frac{4\pi\eta_{\text{th}}k_B}{\hbar}\right)
t_{\text{eff}}\,\frac{T-T_{\text{ref}}}{T_{\text{ref}}}.
\end{equation}

%-------------------------------------------------------------------
\section{Operational Bridge: $I_{\text{path}}$ to Photon Number}
%-------------------------------------------------------------------
From the EGY distinguishability relation~\cite{englert1996}, with coherent
detector states giving $|\langle\alpha_1|\alpha_2\rangle|^2=e^{-\kappa\bar n}$:
\begin{equation}
I_{\text{path}}(\bar n)=\sqrt{1-e^{-\kappa\bar n}}.
\label{eq:n_map}
\end{equation}

%-------------------------------------------------------------------
\section{Statistical Identifiability: AICc Threshold}
%-------------------------------------------------------------------
\begin{equation}
\mathrm{AICc}=N\ln\!\left(\frac{\mathrm{RSS}}{N}\right)
+2K+\frac{2K(K+1)}{N-K-1}.
\end{equation}
\begin{table}[h]
\caption{AICc selection probability ($c=0.5$, $\sigma=0.005$).}
\label{tab:aicc}
\begin{ruledtabular}
\begin{tabular}{cccc}
$N$ & v17.0-C (\%) & v18.0-C (\%) & Reliability\\\hline
8 & 89.0 & 11.0 & Low\\
16 & 57.0 & 43.0 & Low\\
36 & 21.0 & 79.0 & Moderate--High\\
50 & 10.0 & 90.0 & High\\
80 & 5.0 & 95.0 & Very High\\
\end{tabular}
\end{ruledtabular}
\end{table}
\begin{figure}[t]
\centering
\includegraphics[width=0.85\linewidth]{figures/figure2.png}
\caption{AICc selection probability vs.\ $N$.
Threshold: $N\ge36$.}
\label{fig:aicc_threshold}
\end{figure}

%-------------------------------------------------------------------
\section{Unified DTQEM Framework (v63.2-C)}
%-------------------------------------------------------------------
\subsection{Scaling Exponents}
\begin{enumerate}
\item $\beta=0.44$: Debye--Pikovski frozen modes, $\tau_c\propto m^{-0.44}$.
\item $\delta=1/3$: van~der~Waals scattering (eikonal), $\tau_c\propto v^{-1/3}$.
\item $\zeta=0.005$: symmetry-suppressed blackbody emission.
\end{enumerate}
\subsection{Unified Coherence Time}
$C_{\text{joint}}$ is the cross-correlation strength between the kinetic and
thermal baths:
\begin{widetext}
\begin{equation}
\tau_c=\frac{\tau_{c0}}{
A\!\left(\frac{m}{m_u}\right)^{\!\beta}\!\left(\frac{v}{c}\right)^{\!\delta}
I_{\text{path}}
+B(1+\zeta N)\frac{\Delta T}{T_{\text{ref}}}
+C_{\text{joint}}\!\left(\frac{m}{m_u}\right)^{\!\beta}
\!\left(\frac{v}{c}\right)^{\!\delta}
I_{\text{path}}(1+\zeta N)\frac{\Delta T}{T_{\text{ref}}}}.
\label{eq:unified}
\end{equation}
\end{widetext}
\begin{figure*}[t]
\centering
\includegraphics[width=\textwidth]{figures/figure3.png}
\caption{$\tau_c$ landscape: (a)~C$_{60}$, (b)~C$_{700}$.}
\label{fig:unified_landscape}
\end{figure*}

%-------------------------------------------------------------------
\section{Numerical Implementation}
%-------------------------------------------------------------------
All figures and tables were produced by open-source Python scripts archived at
Zenodo~\cite{zenodo_dtqem} and on GitHub
(\url{https://github.com/reddoma742/Davisson-Germer-DTQEM}),
released under CC~BY-NC-SA~4.0.

\begin{table}[h]
\caption{Script-to-output mapping.}
\label{tab:code}
\begin{ruledtabular}
\begin{tabular}{lp{4.8cm}}
Script & Output\\\hline
\texttt{models/v18/dtqem\_joint\_bath\_v18.py}
& Parameter calibration (Tables~\ref{tab:parameters},\,\ref{tab:performance})\\
\texttt{models/unified/dtqem\_unified\_simulator.py}
& AICc analysis (Table~\ref{tab:aicc}, Fig.~\ref{fig:aicc_threshold})\\
\texttt{models/v63/DTQEM\_v63\_1\_C.py}
& $\tau_c$ landscape (Fig.~\ref{fig:unified_landscape})\\
\texttt{models/v16/dtqem\_tunneling\_v16\_1\_C.py}
& Zeno suppression (Fig.~\ref{fig:zeno})\\
\texttt{scripts/generate\_esd\_final.py}
& ESD landscape (Fig.~\ref{fig:esd_landscape})\\
\texttt{dtqem\_chemistry\_fit.py}
& Physical-chemistry bridge (Fig.~\ref{fig:chem_fit})\\
\end{tabular}
\end{ruledtabular}
\end{table}

A related script, \texttt{dtqem\_chemistry\_fit.py}, explores the extension of
DTQEM to physical chemistry. It demonstrates an inverse relationship between
the DTQEM thermal coefficient $b$ and the bond force constant $k$ for a series
of polar hydride molecules, achieving an $R^2\approx0.85$. The code and its
output figure (\texttt{figure\_chemistry\_fit.png}) are archived with the
project, indicating a promising interdisciplinary bridge between quantum
decoherence theory and molecular bonding.

A minimal example of the core class is in Appendix~\ref{app:code}; the full
implementation is at~\cite{zenodo_dtqem}.

%-------------------------------------------------------------------
\section{Entanglement Sudden Death}
%-------------------------------------------------------------------
For a two-qubit X-state under local DTQEM dephasing:
\begin{equation}
\Gamma=\alpha\,I_{\text{path}}+\beta\,\frac{\Delta T}{T_{\text{ref}}}
+\gamma\,I_{\text{path}}\,\frac{\Delta T}{T_{\text{ref}}},
\label{eq:gamma_esd}
\end{equation}
\begin{equation}
\mathcal C(t)=2\max\!\left\{0,\,
|\rho_{14}(0)|\,e^{-2\Gamma t}-\sqrt{\rho_{22}(0)\rho_{33}(0)}\right\},
\end{equation}
\begin{equation}
t_{\text{ESD}}=\frac{K}{\Gamma},\quad
K=\tfrac12\ln\!\left(\frac{|\rho_{14}(0)|}
{\sqrt{\rho_{22}(0)\rho_{33}(0)}}\right).
\label{eq:esd_main}
\end{equation}
\begin{figure}[t]
\centering
\includegraphics[width=0.85\linewidth]{figures/figure_esd_final.png}
\caption{ESD landscape ($\alpha=1.0$, $\beta=0.85$, $\gamma=0.5$).}
\label{fig:esd_landscape}
\end{figure}
\subsection{Quantum Zeno Effect}
\begin{figure}[t]
\centering
\includegraphics[width=0.85\linewidth]{figures/figure_zeno.png}
\caption{Zeno tunneling suppression under continuous measurement.}
\label{fig:zeno}
\end{figure}

%-------------------------------------------------------------------
\section{Discussion}
%-------------------------------------------------------------------
\subsection{Conceptual Note on Quantum Structure}
The hydrogen atom is best described not as two separate bodies, but as a
single quantum system in which the electron and nucleus orbit a common center
of mass, with the finite nuclear mass accounted for through the reduced mass
$\mu_r = m_e m_p/(m_e+m_p)$~\cite{weiss2012}. Consequently, a quantum jump
is not a mechanical attempt to maintain a constant velocity ratio between the
two particles, but rather a transition between discrete, quantized energy
eigenstates. The atomic emission spectrum is therefore a direct signature of
this quantized energy structure---the spectral lines encode the allowed
energy differences---rather than a trace of classical orbital motion. This
perspective reinforces the DTQEM picture: coherence loss reflects the
departure from a pure quantum eigenstate, not a disruption of classical
trajectories.

\subsection{Physical Significance of the Conceptual Note}
Some descriptions in this work use intuitive language to aid physical
interpretation. Such formulations complement rather than replace the
mathematical statements, and were refined through iterative cross-checking.

\subsection{Limitations}
The model assumes pure dephasing ($T_1\gg T_\phi$) and Markovian dynamics.
At ultra-short times or strong coupling, non-Markovian effects appear,
producing coherence revivals~\cite{breuer2016,rivas2014}.

\subsection{Comparison with Experiments}
DTQEM results show strong qualitative agreement with superconducting-qubit
dispersive readout and matter-wave interference of large molecules.
The present comparisons are simulation-based benchmarks; full quantitative
validation against dedicated laboratory datasets remains future work.

\subsection{Experimental Proposals}
\begin{itemize}
\item \textbf{Molecular interferometry:} Talbot--Lau with
$\mathrm{C}_{60}$/$\mathrm{C}_{70}$/$\mathrm{C}_{84}$.
\item \textbf{Superconducting qubits:} Transmon in a readout resonator;
$I_{\text{path}}\leftrightarrow\bar n$ via dispersive coupling;
$T$ tuned in a dilution refrigerator.
\end{itemize}

%-------------------------------------------------------------------
\section{Conclusions}
%-------------------------------------------------------------------
We have presented DTQEM v18.0-C with (i)~first-principles derivations of
$a(\omega_c,T)$ and $b(T)$, (ii)~the operational mapping
$I_{\text{path}}(\bar n)$, (iii)~AICc threshold $N\ge36$,
(iv)~unified v63.2-C scaling, (v)~closed-form ESD predictions, and
(vi)~a Morse-potential bridge linking $b$ to molecular bond force constants.
All results are reproducible via the archived code~\cite{zenodo_dtqem},
providing a practical roadmap for experimental verification in
molecular interferometry and solid-state platforms.

\subsection*{Outlook}
The present work constitutes version~1.0 of the DTQEM interdisciplinary
framework. Future directions include: (a)~expanding the physical-chemistry
database to at least 30 molecules spanning both polar and non-polar bonds, to
test the universality of the $b(k)$ relation; (b)~replacing the local
approximation $b\propto1/k$ with the full Morse-Arrhenius expression
$b(k)=\tau_{\text{eff}}(2\pi)^{-1}\sqrt{k/\mu}\exp[-k/(2a^2k_BT_{\text{ref}})]$
fitted to experimental vibrational data; (c)~extending the AICc threshold
analysis to the chemistry dataset to determine the minimum number of molecules
required for statistical identification of the $b$--$k$ coupling; and
(d)~exploring the non-Markovian corrections relevant at ultrashort times and
strong system-bath coupling.

\begin{acknowledgments}
AI language models were used as analytical aids:
\textbf{DeepSeek} (critical analysis, methodology),
\textbf{Claude/Anthropic} (code, derivations, Zeno, physical-chemistry bridge),
\textbf{Arena~AI} (first-principles derivations, ESD, Morse derivation).
All AI assistance operated under full human supervision.
Scientific claims and validation are solely the responsibility of
\textbf{Reddouane Berramdane}.
\end{acknowledgments}

%-------------------------------------------------------------------
\appendix
%-------------------------------------------------------------------
\section{Core Model Code}
\label{app:code}
\begin{lstlisting}[language=Python]
import numpy as np

class DTQEMModel:
    def __init__(self, C0=0.3675, a=1.6968,
                 b=0.8055, c=0.5000, T_ref=300.0):
        self.C0, self.a, self.b = C0, a, b
        self.c, self.T_ref = c, T_ref

    @staticmethod
    def Ipath_from_nbar(n_bar, kappa=1.0):
        return np.sqrt(1.0 - np.exp(-kappa*np.maximum(0.0, n_bar)))

    def predict(self, I_path, T):
        I = np.clip(I_path, 0.0, 1.0)
        dT = (np.asarray(T) - self.T_ref) / self.T_ref
        return self.C0 * np.exp(
            -self.a*I - self.b*dT - self.c*I*dT)
\end{lstlisting}
Full code: \url{https://github.com/reddoma742/Davisson-Germer-DTQEM} \\
DOI: \href{https://doi.org/10.5281/zenodo.20516742}{10.5281/zenodo.20516742}
(CC~BY-NC-SA~4.0).

%-------------------------------------------------------------------
\section{Theoretical Derivation of the Thermal Coefficient $b$ from the Morse Potential}
\label{app:morse}
%-------------------------------------------------------------------

\subsection{Morse Potential and Force Constant}
The interatomic potential of a diatomic bond is described by the Morse
potential~\cite{morse1929}:
\begin{equation}
V(r)=D_e\!\left[1-e^{-\alpha(r-r_0)}\right]^2,
\label{eq:morse}
\end{equation}
where $D_e$ is the well depth (bond dissociation energy), $r_0$ the
equilibrium bond length, and $\alpha$ the range parameter governing well
width. Expanding to second order about $r_0$ yields the harmonic force
constant:
\begin{equation}
k \equiv \left.\frac{d^2V}{dr^2}\right|_{r=r_0} = 2\alpha^2 D_e
\quad\Longrightarrow\quad
D_e = \frac{k}{2\alpha^2}.
\label{eq:Dk}
\end{equation}

\subsection{Thermally Activated Dissociation Rate}
According to transition-state theory~\cite{eyring1935}, the dissociation
rate at temperature $T$ is
\begin{equation}
\Gamma_{\text{diss}}(T)=\nu_0\exp\!\left(-\frac{D_e}{k_BT}\right),
\label{eq:arrhenius}
\end{equation}
where the attempt frequency $\nu_0$ is the zero-point vibrational frequency:
\begin{equation}
\nu_0=\frac{1}{2\pi}\sqrt{\frac{k}{\mu}},
\label{eq:nu0}
\end{equation}
with $\mu$ the reduced mass of the bond.

\subsection{Connection to the DTQEM Thermal Coefficient}
We identify the DTQEM thermal coefficient $b$ at the reference temperature
$T_{\text{ref}}$ as the product of the effective interaction time
$\tau_{\text{eff}}$ and the thermally activated dissociation rate:
\begin{equation}
b \equiv \tau_{\text{eff}}\,\Gamma_{\text{diss}}(T_{\text{ref}}).
\label{eq:b_def}
\end{equation}
Substituting Eqs.~(\ref{eq:Dk})--(\ref{eq:nu0}) into Eq.~(\ref{eq:b_def})
gives the central result of this appendix:
\begin{equation}
\boxed{
b(k)=\frac{\tau_{\text{eff}}}{2\pi}\sqrt{\frac{k}{\mu}}
\exp\!\left(-\frac{k}{2\alpha^2 k_BT_{\text{ref}}}\right).
}
\label{eq:b_morse}
\end{equation}

\subsection{Physical Interpretation}
Equation~(\ref{eq:b_morse}) encodes a competition between two opposing
effects of bond stiffness:
\begin{itemize}
\item The prefactor $\sqrt{k/\mu}$ increases with $k$, reflecting the higher
      attempt frequency of stiffer bonds.
\item The exponential factor $\exp[-k/(2\alpha^2k_BT_{\text{ref}})]$
      decreases sharply with $k$, because stiffer bonds have deeper wells
      ($D_e=k/2\alpha^2$) and are exponentially harder to dissociate
      thermally.
\end{itemize}
In the experimentally relevant regime of polar hydride molecules
($k\gg k_{\text{th}}\equiv 2\alpha^2 k_BT_{\text{ref}}$), the exponential
dominates and the effective local approximation $b\propto 1/k$ holds.
This is \emph{not} an independent empirical law but a local limit of
Eq.~(\ref{eq:b_morse}) in the strong-bond regime.

\subsection{Preliminary Experimental Validation}
A Python implementation (\texttt{dtqem\_chemistry\_fit.py}) tests
Eq.~(\ref{eq:b_morse}) on five polar hydride molecules
(HF, H$_2$O, HCl, H$_2$S, HBr). A linear fit between the force constant
$k$ and the experimental thermal dissociation temperature $T_{\text{diss}}$
yields $R^2\approx0.85$ ($p=0.026$), consistent with the predicted monotonic
increase of $D_e$ with $k$. The derived $b$ values follow the expected
inverse trend. The corresponding figure is shown in Fig.~\ref{fig:chem_fit}.

\begin{figure}[h]
\centering
\includegraphics[width=0.85\linewidth]{figures/figure_chemistry_fit.png}
\caption{DTQEM interdisciplinary bridge: bond force constant $k$ vs.\
thermal dissociation temperature $T_{\text{diss}}$ (blue, left axis) and
derived DTQEM coefficient $b$ (red, right axis) for five polar hydride
molecules. The dashed red curve shows $b\propto 1/k$ as the local
strong-bond approximation to Eq.~(\ref{eq:b_morse}). Linear fit:
$R^2\approx0.85$.}
\label{fig:chem_fit}
\end{figure}

These results are preliminary; a statistically robust validation requires
at least 30 molecules spanning both polar and non-polar bonds (see Outlook).

%-------------------------------------------------------------------
\begin{thebibliography}{99}
\bibitem{breuer2002}
H.-P.~Breuer and F.~Petruccione,
\emph{The Theory of Open Quantum Systems} (OUP, 2002).
\bibitem{englert1996}
B.-G.~Englert, \emph{Phys.\ Rev.\ Lett.}\ \textbf{77}, 2154 (1996).
\bibitem{weiss2012}
U.~Weiss, \emph{Quantum Dissipative Systems}, 4th ed.\ (World Scientific, 2012).
\bibitem{breuer2016}
H.-P.~Breuer \emph{et al.},
\emph{Rev.\ Mod.\ Phys.}\ \textbf{88}, 021002 (2016).
\bibitem{rivas2014}
A.~Rivas, S.\,F.~Huelga, and M.\,B.~Plenio,
\emph{Rep.\ Prog.\ Phys.}\ \textbf{77}, 094001 (2014).
\bibitem{burnham2002}
K.\,P.~Burnham and D.\,R.~Anderson,
\emph{Model Selection and Multimodel Inference}, 2nd ed.\ (Springer, 2002).
\bibitem{yu2009sudden}
T.~Yu and J.\,H.~Eberly, \emph{Science}\ \textbf{323}, 598 (2009).
\bibitem{morse1929}
P.\,M.~Morse, \emph{Phys.\ Rev.}\ \textbf{34}, 57 (1929).
\bibitem{eyring1935}
H.~Eyring, \emph{J.\ Chem.\ Phys.}\ \textbf{3}, 107 (1935).
\bibitem{zenodo_dtqem}
R.~Berramdane,
\emph{DTQEM v18.0-C: Unified Quantum Decoherence Framework},
Zenodo (2026),
DOI:\,\href{https://doi.org/10.5281/zenodo.20516742}{10.5281/zenodo.20516742}.
\end{thebibliography}

\end{document}
